% ================= SESSION 9.1 =================
\section{9.1 Q-테이블의 한계와 함수근사}
\begin{frame}{Q-테이블이 감당 못 하는 규모 --- 숫자로 셈한다}
  \begin{itemize}
    \item Week 6 Q-learning: $Q[s][a]$ \kb{표}에 모든 상태-행동 쌍 저장.
    \item Atari 화면이 만들 수 있는 상태는 사실상 무한 --- 담을 수도,
      방문하기도 어렵다.
    \item DQN 논문이 \textbf{실제로} 쓴 화면: $84\times84$ \textbf{흑백},
      픽셀당 256가지 밝기.
  \end{itemize}
  \vskip0.6em
  \[
  256^{84 \times 84} = 256^{7{,}056} \approx 10^{16{,}993}
  \]
  \begin{itemize}
    \item 바둑판 $10^{172}$ 보다 \textbf{$10^{16{,}800}$배 더 큰} 상태 공간.
    \item 관측 가능 우주의 원자 수($10^{80}$)는 말할 나위도 없다.
  \end{itemize}
  \vskip0.4em
  \kq{Q-테이블의 문제는 ``메모리가 부족''이 아니라 \textbf{원칙적으로} 불가능하다.}
\end{frame}
